% Options for packages loaded elsewhere
\PassOptionsToPackage{unicode}{hyperref}
\PassOptionsToPackage{hyphens}{url}
\documentclass[
]{article}
\usepackage{xcolor}
\usepackage[margin=1in]{geometry}
\usepackage{amsmath,amssymb}
\setcounter{secnumdepth}{-\maxdimen} % remove section numbering
\usepackage{iftex}
\ifPDFTeX
  \usepackage[T1]{fontenc}
  \usepackage[utf8]{inputenc}
  \usepackage{textcomp} % provide euro and other symbols
\else % if luatex or xetex
  \usepackage{unicode-math} % this also loads fontspec
  \defaultfontfeatures{Scale=MatchLowercase}
  \defaultfontfeatures[\rmfamily]{Ligatures=TeX,Scale=1}
\fi
\usepackage{lmodern}
\ifPDFTeX\else
  % xetex/luatex font selection
\fi
% Use upquote if available, for straight quotes in verbatim environments
\IfFileExists{upquote.sty}{\usepackage{upquote}}{}
\IfFileExists{microtype.sty}{% use microtype if available
  \usepackage[]{microtype}
  \UseMicrotypeSet[protrusion]{basicmath} % disable protrusion for tt fonts
}{}
\makeatletter
\@ifundefined{KOMAClassName}{% if non-KOMA class
  \IfFileExists{parskip.sty}{%
    \usepackage{parskip}
  }{% else
    \setlength{\parindent}{0pt}
    \setlength{\parskip}{6pt plus 2pt minus 1pt}}
}{% if KOMA class
  \KOMAoptions{parskip=half}}
\makeatother
\usepackage{color}
\usepackage{fancyvrb}
\newcommand{\VerbBar}{|}
\newcommand{\VERB}{\Verb[commandchars=\\\{\}]}
\DefineVerbatimEnvironment{Highlighting}{Verbatim}{commandchars=\\\{\}}
% Add ',fontsize=\small' for more characters per line
\usepackage{framed}
\definecolor{shadecolor}{RGB}{248,248,248}
\newenvironment{Shaded}{\begin{snugshade}}{\end{snugshade}}
\newcommand{\AlertTok}[1]{\textcolor[rgb]{0.94,0.16,0.16}{#1}}
\newcommand{\AnnotationTok}[1]{\textcolor[rgb]{0.56,0.35,0.01}{\textbf{\textit{#1}}}}
\newcommand{\AttributeTok}[1]{\textcolor[rgb]{0.13,0.29,0.53}{#1}}
\newcommand{\BaseNTok}[1]{\textcolor[rgb]{0.00,0.00,0.81}{#1}}
\newcommand{\BuiltInTok}[1]{#1}
\newcommand{\CharTok}[1]{\textcolor[rgb]{0.31,0.60,0.02}{#1}}
\newcommand{\CommentTok}[1]{\textcolor[rgb]{0.56,0.35,0.01}{\textit{#1}}}
\newcommand{\CommentVarTok}[1]{\textcolor[rgb]{0.56,0.35,0.01}{\textbf{\textit{#1}}}}
\newcommand{\ConstantTok}[1]{\textcolor[rgb]{0.56,0.35,0.01}{#1}}
\newcommand{\ControlFlowTok}[1]{\textcolor[rgb]{0.13,0.29,0.53}{\textbf{#1}}}
\newcommand{\DataTypeTok}[1]{\textcolor[rgb]{0.13,0.29,0.53}{#1}}
\newcommand{\DecValTok}[1]{\textcolor[rgb]{0.00,0.00,0.81}{#1}}
\newcommand{\DocumentationTok}[1]{\textcolor[rgb]{0.56,0.35,0.01}{\textbf{\textit{#1}}}}
\newcommand{\ErrorTok}[1]{\textcolor[rgb]{0.64,0.00,0.00}{\textbf{#1}}}
\newcommand{\ExtensionTok}[1]{#1}
\newcommand{\FloatTok}[1]{\textcolor[rgb]{0.00,0.00,0.81}{#1}}
\newcommand{\FunctionTok}[1]{\textcolor[rgb]{0.13,0.29,0.53}{\textbf{#1}}}
\newcommand{\ImportTok}[1]{#1}
\newcommand{\InformationTok}[1]{\textcolor[rgb]{0.56,0.35,0.01}{\textbf{\textit{#1}}}}
\newcommand{\KeywordTok}[1]{\textcolor[rgb]{0.13,0.29,0.53}{\textbf{#1}}}
\newcommand{\NormalTok}[1]{#1}
\newcommand{\OperatorTok}[1]{\textcolor[rgb]{0.81,0.36,0.00}{\textbf{#1}}}
\newcommand{\OtherTok}[1]{\textcolor[rgb]{0.56,0.35,0.01}{#1}}
\newcommand{\PreprocessorTok}[1]{\textcolor[rgb]{0.56,0.35,0.01}{\textit{#1}}}
\newcommand{\RegionMarkerTok}[1]{#1}
\newcommand{\SpecialCharTok}[1]{\textcolor[rgb]{0.81,0.36,0.00}{\textbf{#1}}}
\newcommand{\SpecialStringTok}[1]{\textcolor[rgb]{0.31,0.60,0.02}{#1}}
\newcommand{\StringTok}[1]{\textcolor[rgb]{0.31,0.60,0.02}{#1}}
\newcommand{\VariableTok}[1]{\textcolor[rgb]{0.00,0.00,0.00}{#1}}
\newcommand{\VerbatimStringTok}[1]{\textcolor[rgb]{0.31,0.60,0.02}{#1}}
\newcommand{\WarningTok}[1]{\textcolor[rgb]{0.56,0.35,0.01}{\textbf{\textit{#1}}}}
\usepackage{graphicx}
\makeatletter
\newsavebox\pandoc@box
\newcommand*\pandocbounded[1]{% scales image to fit in text height/width
  \sbox\pandoc@box{#1}%
  \Gscale@div\@tempa{\textheight}{\dimexpr\ht\pandoc@box+\dp\pandoc@box\relax}%
  \Gscale@div\@tempb{\linewidth}{\wd\pandoc@box}%
  \ifdim\@tempb\p@<\@tempa\p@\let\@tempa\@tempb\fi% select the smaller of both
  \ifdim\@tempa\p@<\p@\scalebox{\@tempa}{\usebox\pandoc@box}%
  \else\usebox{\pandoc@box}%
  \fi%
}
% Set default figure placement to htbp
\def\fps@figure{htbp}
\makeatother
\setlength{\emergencystretch}{3em} % prevent overfull lines
\providecommand{\tightlist}{%
  \setlength{\itemsep}{0pt}\setlength{\parskip}{0pt}}
\usepackage{bookmark}
\IfFileExists{xurl.sty}{\usepackage{xurl}}{} % add URL line breaks if available
\urlstyle{same}
\hypersetup{
  pdftitle={The Beta-Binomial Model},
  hidelinks,
  pdfcreator={LaTeX via pandoc}}

\title{The Beta-Binomial Model}
\author{}
\date{\vspace{-2.5em}2026-04-17}

\begin{document}
\maketitle

\subsection{Introduction}\label{introduction}

The beta-binomial conjugate pair is the textbook entry point to Bayesian
inference: a Beta prior on a probability combined with a binomial
likelihood yields a Beta posterior whose parameters are obtained by
adding successes to one shape parameter and failures to the other. The
closed form makes the entire workflow --- prior elicitation, posterior,
credible interval, posterior predictive --- tractable on a calculator.
It is also the conceptual backbone for hierarchical models of
proportions, which use a Beta prior at the population level to share
strength across many groups.

\subsection{Prerequisites}\label{prerequisites}

Familiarity with the Beta distribution as a flexible density on
\([0, 1]\), the binomial likelihood for independent Bernoulli trials,
and the basic mechanics of Bayesian updating.

\subsection{Theory}\label{theory}

With prior \(\theta \sim \mathrm{Beta}(a, b)\) and likelihood
\(X \sim \mathrm{Binomial}(n, \theta)\), the posterior is

\[\theta \mid x \sim \mathrm{Beta}(a + x, b + n - x).\]

The shape parameters of the prior have a direct interpretation as
pseudo-counts: \(a\) pseudo-successes and \(b\) pseudo-failures observed
before the data arrived, so the prior effective sample size is
\(a + b\). The posterior mean is

\[\mathbb E[\theta \mid x] = \frac{a + x}{a + b + n},\]

a weighted average of the prior mean \(a / (a+b)\) and the sample
proportion \(x / n\), with weights determined by the prior pseudo-count
and the actual sample size. Conjugacy also gives closed-form expressions
for the posterior variance, credible intervals, and the predictive
distribution for future Bernoulli trials.

\subsection{Assumptions}\label{assumptions}

Independent Bernoulli trials with a constant success probability
\(\theta\), a Beta prior plausible enough that its tails do not exclude
reasonable values, and a sample size that is fixed in advance ---
informative censoring or stopping rules invalidate the binomial
likelihood.

\subsection{R Implementation}\label{r-implementation}

\begin{Shaded}
\begin{Highlighting}[]
\CommentTok{\# Prior Beta(2, 2) {-}{-} weakly informative around 0.5}
\CommentTok{\# Data: 15 successes in 25 trials}
\NormalTok{a }\OtherTok{\textless{}{-}} \DecValTok{2}\NormalTok{; b }\OtherTok{\textless{}{-}} \DecValTok{2}
\NormalTok{x }\OtherTok{\textless{}{-}} \DecValTok{15}\NormalTok{; n }\OtherTok{\textless{}{-}} \DecValTok{25}

\NormalTok{post\_a }\OtherTok{\textless{}{-}}\NormalTok{ a }\SpecialCharTok{+}\NormalTok{ x}
\NormalTok{post\_b }\OtherTok{\textless{}{-}}\NormalTok{ b }\SpecialCharTok{+}\NormalTok{ n }\SpecialCharTok{{-}}\NormalTok{ x}

\FunctionTok{c}\NormalTok{(}\AttributeTok{prior\_mean =}\NormalTok{ a }\SpecialCharTok{/}\NormalTok{ (a }\SpecialCharTok{+}\NormalTok{ b),}
  \AttributeTok{sample\_prop =}\NormalTok{ x }\SpecialCharTok{/}\NormalTok{ n,}
  \AttributeTok{posterior\_mean =}\NormalTok{ post\_a }\SpecialCharTok{/}\NormalTok{ (post\_a }\SpecialCharTok{+}\NormalTok{ post\_b),}
  \AttributeTok{ci95 =} \FunctionTok{qbeta}\NormalTok{(}\FunctionTok{c}\NormalTok{(}\FloatTok{0.025}\NormalTok{, }\FloatTok{0.975}\NormalTok{), post\_a, post\_b))}

\NormalTok{theta }\OtherTok{\textless{}{-}} \FunctionTok{seq}\NormalTok{(}\DecValTok{0}\NormalTok{, }\DecValTok{1}\NormalTok{, }\AttributeTok{length.out =} \DecValTok{400}\NormalTok{)}
\FunctionTok{plot}\NormalTok{(theta, }\FunctionTok{dbeta}\NormalTok{(theta, a, b), }\AttributeTok{type =} \StringTok{"l"}\NormalTok{, }\AttributeTok{col =} \StringTok{"grey"}\NormalTok{,}
     \AttributeTok{ylab =} \StringTok{"Density"}\NormalTok{, }\AttributeTok{main =} \StringTok{"Beta{-}Binomial update"}\NormalTok{)}
\FunctionTok{lines}\NormalTok{(theta, }\FunctionTok{dbeta}\NormalTok{(theta, post\_a, post\_b), }\AttributeTok{col =} \StringTok{"steelblue"}\NormalTok{, }\AttributeTok{lwd =} \DecValTok{2}\NormalTok{)}
\end{Highlighting}
\end{Shaded}

\subsection{Output \& Results}\label{output-results}

A short numeric vector with the prior mean (0.50), the sample proportion
(0.60), the posterior mean (0.59), and the 95 \% equal-tailed credible
interval (0.40 to 0.76). The prior-versus-posterior plot makes the pull
of the prior visible: with only 25 trials the posterior shifts
noticeably toward the data but retains the curvature of the prior.

\subsection{Interpretation}\label{interpretation}

A reporting sentence: ``A Beta(2, 2) prior combined with 15 successes in
25 trials yielded a Beta(17, 12) posterior with mean 0.59 (95 \%
credible interval 0.40 to 0.76); the data have moved the posterior
\(0.09\) above the prior mean while excluding values below 0.40 with
97.5 \% probability.'' Always disclose the prior in the methods section
so that readers can see how much shrinkage it imposes.

\subsection{Practical Tips}\label{practical-tips}

\begin{itemize}
\tightlist
\item
  \(\mathrm{Beta}(1, 1)\) is the uniform reference prior;
  \(\mathrm{Beta}(0.5, 0.5)\) is Jeffreys; both keep posterior inference
  data-driven.
\item
  For genuinely informative priors, set \(a + b\) equal to the prior
  effective sample size you can defend (e.g.~results from a prior trial
  of size 20 with proportion 0.7 give \(\mathrm{Beta}(14, 6)\)).
\item
  Use \texttt{qbeta()} for credible intervals and \texttt{pbeta()} for
  posterior tail probabilities; HPD intervals come from
  \texttt{HDInterval::hdi(qbeta(p,\ post\_a,\ post\_b))} evaluated on a
  grid.
\item
  For multi-category outcomes, generalise to the Dirichlet--multinomial,
  which inherits the conjugate structure on a probability simplex.
\item
  When the prior is contested, report a sensitivity analysis across at
  least \(\mathrm{Beta}(1, 1)\), \(\mathrm{Beta}(2, 2)\), and one
  informative pair; if the conclusion holds across all three, the result
  is robust to prior choice.
\end{itemize}

\subsection{R Packages Used}\label{r-packages-used}

Base R suffices for the conjugate update (\texttt{dbeta},
\texttt{pbeta}, \texttt{qbeta}, \texttt{rbeta});
\texttt{bayestestR::point\_estimate()} and \texttt{HDInterval::hdi()}
add tidy summaries and HPD intervals when a more polished workflow is
needed.

\subsection{For Reviewers}\label{for-reviewers}

\textbf{What to look for in a paper using this method.}

\begin{itemize}
\tightlist
\item
  \textbf{Common misapplications.}
\item
  \textbf{Diagnostics that should be reported but often aren't.}
\item
  \textbf{Red flags in tables and figures.}
\item
  \textbf{What to verify.}
\item
  \textbf{What an adequate Methods paragraph must contain.}
\end{itemize}

\subsection{See also --- labs in this
chapter}\label{see-also-labs-in-this-chapter}

\begin{itemize}
\tightlist
\item
  \href{labs/lab-bayesian-thinking-control-vs-decision-error.Rmd}{Bayesian
  thinking; α control vs decision error}
\item
  \href{labs/lab-brms-stan-loo-hierarchical-models.Rmd}{brms / Stan,
  LOO, hierarchical models}
\end{itemize}

\end{document}
